\documentclass[11pt]{article}

\usepackage[french]{babel}

\usepackage[T1]{fontenc}

\usepackage{amsfonts}
\usepackage{amsmath}
\usepackage{amssymb}

\usepackage{graphicx}
\usepackage{multirow}
\usepackage{multicol}
\usepackage{xcolor}
\usepackage[shortlabels]{enumitem}

\usepackage[margin=2cm]{geometry}

\usepackage{mathtools}
\usepackage{siunitx}
\usepackage{physics}

\usepackage{tikz}
\usepackage{hyperref}

\usepackage{algorithm}
\usepackage{algpseudocode}
\usepackage{listings}
\usetikzlibrary{quantikz2,babel}

\tikzset{qcircuit/.style={
    column sep=.9cm,
    row sep=0.7cm,
    operator/.style={draw=black!80, line width=0.5pt, fill=gray!5, minimum height=0.7cm, minimum width=0.7cm},
    phase/.style={draw=black!80, fill=black!20, circle, minimum size=4pt},
    slice/.style={font=\footnotesize\itshape, color=gray!50, dashed}
}}


\DeclarePairedDelimiter\ceil{\lceil}{\rceil}
\DeclarePairedDelimiter\floor{\lfloor}{\rfloor}


\AtBeginDocument{\RenewCommandCopy\qty\SI}

\makeatletter
   \def\vhrulefill#1{\leavevmode\leaders\hrule\@height#1\hfill \kern\z@}
\makeatother

\newcommand{\moy}[1]{\langle #1 \rangle}
\newcommand{\pDer}[2]{\frac{\partial #1}{\partial #2}}
\newcommand{\pDerd}[2]{\frac{\partial^2 #1}{\partial #2^2}}

\newcounter{questioncounter}
\newenvironment{question}[1]{
\refstepcounter{questioncounter}
\noindent\rule{0.05\textwidth}{1pt}\,\,\raisebox{-.3\ht\strutbox}{\textbf{Question \thequestioncounter \hspace{5mm} #1}}\,\,\vhrulefill{1pt}\par\vspace{0.4cm}
}{\par}

%======================================================
%en-tete
%======================================================
\begin{document}
\selectlanguage{french}
\noindent{\Large Anthony Dubeau: Duba0207 } \hfill {\Large \today}\\
{\Large Anthony Wroblewski: Wroa8915}


\par\vspace{5mm}
\begin{center}
   {\Large IFT436 - Algorithmes et structures de données} \\ \vspace{5mm}
   {\LARGE\sffamily Devoir \#2} \\ \vspace{5mm}
\end{center}

%======================================================
%question 1
%======================================================
\begin{question}{Retour sur Dijkstra} % j'ai l'impression d'voir un boucle de trop pour le n^3
\begin{enumerate}[label=(\alph*)]
    \item
\begin{lstlisting}[mathescape=true]
D <- {}
pour chaque t $\in$ T 
    visite <- {t:0}
    dernier <- T
    tant que dernier $\notin$ b
        cout <- $\infty$
        pour chaque u $\in$ visite.clees()
            pour chaque v \in N$_g$(u)
                cout_v = visite[u] + f(u,v)
                si cout_v < cout
                    cout <- cout_v
                    dernier <- v
                    visite[v] <- cout_v
                fin si
            fin pour
        fin pour
    fin tant que
    D[t] = visite[dernier]
\end{lstlisting}
\end{enumerate}
\end{question}


%======================================================
%question 2
%======================================================
\begin{question}{Récurrences récurrentes}

\begin{enumerate}[label=(\alph*)]
\item L'énnoncer nous fourni, $a=5, \;b=4 \; et \; c=1,08$
    $$\Rightarrow f(n)=5f(n/4)+\mathcal{O}({n^{1,08}})$$
    C'est le cas 1 du théorème maitre : $ \log_b{a} > c$
    $$\log_b{a}=\log_4{5} \approx 1,16 > c =1,08$$
    $$  \Rightarrow f(n) \in \mathcal{O}({n^{\log_4{5}}})$$
\item Par inspection : $a=2,\;b=2, \; c=?$. On trouve aussi la récurrence suivante:
    $$f(n)=2f(n/2)+2(\mathcal{O}(n)+\mathcal{O}(n\log(n)))$$
    $$f(n)=2f(n/2)+\mathcal{O}(n\log(n))$$
    $$ \Rightarrow c=1\;et\;k=1$$
    Application du théorème maitre : $\log_a{b}=\log_2{2}=1=c$\vspace{5mm}\\
    C'est le cas 2 du théorème maitre : $k = k+1$
    $$\Rightarrow f(n) \in \mathcal{O}(n\log^{2}{n})$$
\item %pas les technique encore pour faire l'analyse d'un algo diviser avec n-1
     
\end{enumerate}
\end{question}


%======================================================
%question 3
%======================================================
\begin{question}{Divisez et régnez}
    \begin{enumerate}[label=(\alph*)]
\item  
\end{enumerate}
\end{question}

%======================================================
%question 4
%======================================================
\begin{question}{La gestion de conflits de bar}
    \begin{enumerate}[label=(\alph*)]
\item  
\end{enumerate}
\end{question}

\bibliographystyle{apsrev4-2}
\bibliography{references}
\end{document}
